\title{Large Language Models Can Automatically Engineer Features for Few-Shot Tabular Learning}

\begin{abstract}
Large Language Models (LLMs) exhibit exceptional proficiency in handling complex and novel reasoning challenges, presenting significant promise for advancing tabular learning, which underpins a host of practical applications. This work introduces an innovative in-context learning strategy, \textsf{FeatLLM}, in which LLMs serve as feature generators to craft an input dataset tailored for optimal tabular prediction. The features produced through this process allow a straightforward downstream machine learning algorithm—such as linear regression—to estimate class probabilities, thereby facilitating effective few-shot learning. Importantly, \textsf{FeatLLM} relies exclusively on this lightweight predictive model and the extracted features during inference, eliminating the necessity of querying the LLM for every individual sample at test time. Additionally, the approach requires nothing more than API-based interaction with LLMs and navigates issues arising from prompt length constraints. Empirical results spanning multiple tabular benchmarks across various domains highlight that \textsf{FeatLLM} creates high-quality feature rules, achieving a notable average improvement of 10\% over state-of-the-art baselines like TabLLM and STUNT.
\end{abstract}

\section{Introduction}
Large Language Models (LLMs) have recently demonstrated a strong capacity for producing pertinent outputs when faced with tasks outside their training set, negating the need for specialized fine-tuning for new tasks~\cite{creswell2022selection,imani2023mathprompter,taylor2022galactica}. Because these models are trained on massive amounts of text data, they amass a broad reservoir of prior knowledge, which can fill the gap traditionally occupied by human expertise in a range of fields~\cite{Genomes2021,singhal2023large,wilcox2003role}. Such capability is key to automating various processes, including but not limited to dialogue analysis~\cite{finch2023leveraging} and automated program repair~\cite{fan2023automated}. Specifically, by incorporating both the task's description and a small collection of annotated examples within the prompt, LLMs can contextualize the problem, selectively focusing on pertinent features and observations~\cite{brown2020language,zhao2023survey}. This ability suggests that LLMs could operate effectively as feature engineering experts. 

Efforts to harness the generalization and reasoning powers of LLMs are increasingly being extended to tabular data analysis. Domain knowledge embedded in LLMs, such as associations between age and certain medical risks, can inform tasks like disease prediction. In such contexts, prior work encodes task definitions and features into natural language, converts tabular data into serialized formats, and supplies this input to an LLM~\cite{dinh2022lift,wang2023anypredict}. Subsequent modeling steps often utilize in-context learning~\cite{nam2023semi} or apply efficient fine-tuning approaches~\cite{dinh2022lift,hegselmann2023tabllm}. The utilization of LLM-based priors has been demonstrated to deliver robust performance, frequently surpassing classical tabular learning models in scenarios with limited data~\cite{hegselmann2023tabllm}.

Existing tabular learning strategies utilizing large language models (LLMs) encounter several notable challenges. For instance, making end-to-end predictions necessitates at least one LLM invocation per data instance, leading to substantial computational overhead. Achieving optimal accuracy typically involves fine-tuning the LLM, which is impractical for many LLMs lacking publicly available training frameworks or tuning interfaces—a pressing issue as many state-of-the-art LLMs restrict user access to inference APIs only~\cite{anil2023palm,openai2023gpt4}. Moreover, current methodologies struggle with extended prompt lengths~\cite{liu2023lost}; when tabular data contains numerous features, and their serialization results in elongated text prompts, system feasibility diminishes and model performance often suffers. These factors considerably hinder the adoption of LLM-based tabular learning in real-world contexts.

Rather than using LLMs strictly for end-to-end prediction tasks, our method harnesses their capacity for feature engineering. We conceptualize the LLM as a means to leverage its inherent knowledge more effectively, enabling it to generate informative features. Concretely, we propose an innovative in-context learning strategy, \textsf{FeatLLM}, designed to capture the implicit “criteria” behind LLM outputs. For example, in disease prediction scenarios, the LLM can elucidate and articulate the rules or conditions involving features that characterize the target disease. These generated criteria serve as new features, which can replace the original features for task resolution. This approach repositions the LLM's functionality towards feature creation, allowing downstream models to be simplified—once these features are identified, inference can proceed with less complex models, thereby reducing inference time. Exploiting the in-context learning strengths of LLMs, features can be obtained simply by incorporating exemplar cases into prompts, obviating the need for traditional model training. To address the constraint of oversized prompts, ensemble methods~\cite{breiman2001random}, such as feature bagging, can be adopted to manage and minimize prompt length.

\textsf{FeatLLM} organizes the process of generating novel features through prompting into two distinct sub-tasks: (1) comprehending the problem context and deducing how features relate to the target variable; and (2) utilizing insights obtained from this analysis along with a handful of example cases to formulate rules that effectively separate classes. This systematic reasoning method guides LLMs to disregard irrelevant features during rule creation. The formulated rules are executed on the dataset (as parsed program instructions), converting each instance into a binary indicator that reflects compliance with the specific rule. Subsequently, a simple model is trained on these binary features to estimate class probabilities. To enhance reliability, an ensemble mechanism is applied, repeatedly constructing features and employing feature bagging. By carefully tuning the parameters governing the number of features and samples involved in bagging, \textsf{FeatLLM} addresses issues associated with overly lengthy prompts. It circumvents the necessity for model training and incurs minimal inference costs, thereby facilitating practical deployment of LLMs across a range of application domains.

\textsf{FeatLLM} was tested on 13 tabular datasets under low-sample conditions, demonstrating consistently strong and dependable results. Empirical findings indicate that LLMs effectively leverage both existing knowledge and information obtained from the data to extract high-quality features. Our approach surpasses state-of-the-art few-shot learning baselines in diverse experimental configurations. The source code is available at an anonymous GitHub repository:~\texttt{https://github.com/Sungwon-Han/FeatLLM}.

\section{Related Work}

\subsection{Few-Shot Learning with Tabular Data} Advances in few-shot learning have made it possible to derive broadly useful data representations despite limited labeled instances~\cite{chen2018closer}. Although the field initially concentrated on images, few-shot learning frameworks now exhibit strong transferability across various domains, such as text, audio, and increasingly, tabular data~\cite{majumder2022few,nam2023stunt,schick2021exploiting}. Training models with scarce labeled examples, however, increases the susceptibility to erroneous correlations~\cite{bartlett2021deep}. To mitigate this, contemporary research has leveraged auxiliary data sources or incorporated domain-specific priors to embed favorable inductive biases during training. For example, semi-supervised approaches have taken advantage of plentiful unlabeled tabular data, enhanced by strategic augmentation~\cite{ucar2021subtab,yoon2020vime}, or have relied on task-neutral unsupervised pre-training for model initialization~\cite{somepalli2021saint}. Typical unsupervised objectives involve deliberate data masking or corruption and subsequently reconstructing the data~\cite{arik2021tabnet,majmundar2022met}, as well as deploying augmentation strategies for contrastive learning~\cite{bahri2022scarf,chen2020simple}. Nonetheless, due to the inherent heterogeneity of tabular datasets, these augmentation methods are difficult to generalize effectively, and their efficacy has been limited in few-shot scenarios~\cite{nam2023stunt}.

Emerging studies advocate for training on a wide spectrum of tabular datasets, regardless of their relevance to the target problem, followed by transfer learning for downstream tasks~\cite{zhang2023generative,levin2022transfer,zhu2023xtab}. An illustration is TransTab~\cite{wang2022transtab}, where transformer architectures are used to learn the semantic interplay among columns based on their names as well as their values. This strategy has empowered the application of zero-shot and few-shot inference to unseen tabular tasks by utilizing knowledge acquired from diverse tables and column configurations. Conversely, TabPFN~\cite{hollmann2022tabpfn} employs mixtures of distributions to synthesize datasets, pre-training a Bayesian neural network on these. Following pre-training, inference involves supplying both training and test data from the target task, with no additional model training required. The prior knowledge gleaned from various datasets enhances the method’s performance relative to classical tree-based algorithms and has shown notable benefit under few-shot conditions.

\subsection{Language-Interfaced Tabular Learning} Leveraging prior information can also be achieved through the integration of large language models (LLMs)~\cite{anil2023palm,openai2023gpt4}. Trained on vast and varied text datasets, LLMs have demonstrated remarkable generalization to unseen problems and are effective across a breadth of domains~\cite{brown2020language}. Recent work aims to utilize the extensive prior knowledge embedded within LLMs for tabular learning tasks by converting tables into textual format and providing this serialized data as input prompts to the models~\cite{dinh2022lift,hegselmann2023tabllm,wang2023anypredict}. By incorporating feature details, task information, and tabular entries in prompts, LLMs can deduce relationships between variables and tasks to support inference. Contemporary research has pursued two main directions: specifically adapting LLMs to tabular tasks using parameter-efficient tuning methods such as LoRA~\cite{hu2021lora} or IA3~\cite{liu2022few}, or employing in-context learning strategies, wherein several example tasks are included in the prompts to facilitate inference without model updates~\cite{dinh2022lift,hegselmann2023tabllm,nam2023semi,slack2023tablet}.

Although the above approaches demonstrate strong performance in enhancing few-shot tabular learning, the necessity to route inference through an LLM for every instance introduces complications, particularly for deployment in industrial settings. This work proposes moving away from the prevailing black-box paradigm that relies on LLMs to produce end-to-end predictions for each data point. Rather, it leverages the LLM to directly deduce the class-specific conditions or rule sets. \textsf{FeatLLM} takes these inferred rules and translates them into executable code to construct new features, so predictions can occur without LLM inference. By drawing upon in-context learning, the model capitalizes on existing knowledge and limited sample data to generate useful rules.

\section{Methods}
\textsf{FeatLLM} focuses on mining "rules" — logical conditions that define each response class — from the small number of provided examples, as visualized in Figure~\ref{fig:main_model}. These extracted conditions are interpreted by an LLM and subsequently used to generate new binary features for the input sample, each signifying whether a rule is fulfilled or not. The collection of these binary indicators is then fed into a dot product computation with trainable, non-negative weights to approximate the probability for each class. Through model training, the system learns how significant each feature is for predicting class outcomes (see Section~\ref{Sec:training}). To further enhance performance, this rule-based feature extraction and prediction pipeline is repeated with bagging, and results are aggregated by ensembling. The precise methodology and the nuances of each step are detailed below.

\vspace*{-0.12in}
\paragraph{Problem formulation.} Consider a tabular dataset composed of $N$ labeled instances $\mathcal{D} = \{(\mathbf{x}^i, \mathbf{y}^i)\}_{i=1}^{N}$, where each input $\mathbf{x}^i$ is a vector of $d$ features. Each feature is named in natural language, forming the set $F = \{f_j\}_{j=1}^{d}$, and their corresponding definitions are specified separately. The labels $\mathbf{y}^i$ are assumed to come from a predetermined set of classes in a classification framework. For $k$-shot learning, we restrict the training data to $k$ randomly selected labeled instances ($k < N$).

\subsection{Prompt Design for Extracting Rules}
\label{Sec:prompt_design}
In order to facilitate rule extraction with precise reasoning by the LLM, we structure the problem-solving approach to emulate the workflow of expert practitioners dealing with tabular data. Our constructed prompt integrates three principal components, detailed below (refer to Fig.~\ref{fig:prompt_for_rule} and Appendix~\ref{sec:example_rule}):

\vspace*{-0.12in}
\paragraph{Basic information description.} The initial segment delivers core information needed for task resolution (illustrated by the orange annotations in Fig.~\ref{fig:prompt_for_rule}). This section encompasses a statement of the task (including label and feature context) and illustrative examples. The task itself is framed as a question, for example, ``Does this patient's myocardial infarction complications data indicate chronic heart failure? Yes or no?" Additional instances are documented in Appendix~\ref{sec:task_desc}. Feature descriptions specify the data type and may incorporate feature definitions; for categorical attributes, their possible values are enumerated. As part of the demonstration, several training examples (few-shot samples) are rendered as text and paired with their respective correct labels, adhering to the serialization protocols established in earlier work~\cite{dinh2022lift,hegselmann2023tabllm}:

\begin{align}
&\text{Serialize}(\mathbf{x}^i,\mathbf{y}^i, F) = \nonumber \\ 
&``\ f_1\ \text{is}\ \mathbf{x}^i_1.\ ... \ f_d\  \text{is}\  \mathbf{x}^i_d.\ \ \text{Answer:}\  \mathbf{y}^i\ ”
\end{align}

\vspace*{-0.12in}
\paragraph{Reasoning instruction.} Our approach seeks to strengthen the LLM's reasoning capabilities by incorporating explicit reasoning instructions, as illustrated by the blue annotations in Fig.~\ref{fig:prompt_for_rule}. The reasoning guidance commences with an opening statement reminiscent of the chain-of-thought methodology~\cite{wei2022chain}. The inference process is then systematically divided into two distinct phases. First, the LLM is prompted to hypothesize the causal dynamics or tendencies between task-related features and the task description, utilizing its general knowledge or common sense, but without any reference to provided demonstrations. This enables the LLM to utilize its prior understanding of the problem domain. In the subsequent phase, the LLM draws upon both the demonstration examples and insights gleaned in the initial phase to infer decision rules for each class. This bifurcated reasoning strategy mitigates the risk of the model relying on misleading correlations from unimportant features, promoting attention to more pertinent feature attributes. To avoid both the risks of overfitting through excessive rule extraction and underfitting due to insufficient rules, a fixed quota (10 rules)\footnote{Further exploration on rule count is detailed in Section~\ref{sec:analysis}.} is specified for each class. Furthermore, when generating these rules, the LLM references the feature type information from the basic description to prevent rule type inconsistencies.

\vspace*{-0.12in}
\paragraph{Response instruction.} To improve the interpretability and usability of the extracted rules, we direct the LLM via carefully crafted response instructions, shown as yellow text in Fig.~\ref{fig:prompt_for_rule}. The prompt specifies the precise output format expected for each answer class and stipulates the structural template for formulating each rule. This guidance permits the LLM to tailor its responses based on the types of features involved. Even without supplemental instructions, the model is capable of synthesizing multiple rules using logical conjunctions (e.g., AND, OR). Illustrative results are included in Appendix~\ref{sec:outcome-rule}.

\subsection{Inferring Class Likelihood via Rules}
\label{Sec:training}

\paragraph{Rule parsing for feature derivation.} The rules obtained from the earlier stage are employed to produce novel binary features. For each class, these features reflect whether a particular sample fulfills the rule criteria tied to that class, marked by either ‘0’ or ‘1’. Such features facilitate predictions regarding a sample's association with each class. Because the rules are formulated in natural language by the LLM, a tailored parsing method must be developed to automate this feature extraction process. Although the formatting of LLM responses and rule statements is restricted to support easier parsing, outputs may still exhibit varied or inconsistent syntax~\cite{kovavc2023large}. For example, the LLM might replace symbols like ``$>$'' and ``$<$'' with their word equivalents such as ``greater than'' or ``less than,'' which complicates the parsing procedure. 

To overcome issues posed by noisy and inconsistent text parsing, instead of implementing intricate programming logic, we utilize the LLM directly for the task. The LLM demonstrates robust semantic comprehension even in the presence of syntactic variability and is capable of producing concise parsing code based on instructions (thanks to its code-aware training corpus). To capitalize on these strengths, we structure the prompt to include function names, descriptions of input and output, and the extracted rules before submitting it to the LLM. Illustrative prompt examples and their outputs are depicted in Figures~\ref{fig:prompt_for_parsing} and ~\ref{sec:example_parsing}-\ref{sec:example_parse_outcome} in the Appendix. The resulting code snippets are executed via Python’s \texttt{exec()} functionality, enabling efficient data transformation.

\vspace*{-0.12in}
\paragraph{Class likelihood estimation.} With these newly obtained binary rule-based features per class, one straightforward approach to gauge the likelihood that a sample belongs to a given class is to tally the number of class-specific rules the sample satisfies—effectively, the sum of binary features for each class. Yet, recognizing that individual rules and features do not equally influence predictions, it becomes necessary to ascertain their relevance by learning from training data. This is achieved by fitting a standard machine learning (ML) model to the binary features associated with each class, thus estimating the weight or importance of each rule. For illustration, one can apply a bias-free linear model. Let the binary feature vector for sample $\mathbf{x}^i$ regarding class $k$ be defined as $\mathbf{z}_k^i \in \{0, 1\}^R$, where $R$ stands for the total rules in each class and $c$ denotes the number of distinct classes. The class probability vector $\mathbf{p}^i$ is then calculated as follows:

\begin{align}
\text{logit}_k^i &= \max(\mathbf{w}_k, 0) \cdot \mathbf{z}_k^i, \nonumber \\
\mathbf{p}^i &= \text{Softmax}({[\text{logit}_{1}^i, ..., \text{logit}_{c}^i]}). \label{eq:infer_prob}
\end{align}

In this context, $\mathbf{w}_k$ denotes the adjustable parameters of the linear model, which reflect the relative influence of each feature. By restricting these weights to positive values, we guarantee that the features derived from the LLM will not negatively affect the predicted probability of a class during the learning process. Afterward, the logits are transformed into class probabilities using the softmax operation. To update the weights $\mathbf{w}_k$, the cross-entropy loss is minimized over a limited number of training examples.

\vspace*{-0.12in}
\paragraph{Ensembling with bagging.} The method subsequently repeats the entire training and inference cycle multiple times, constructing an ensemble of models for prediction. With $T$ sets of learned weights, $\{\mathbf{w}[t]\}_{t=1}^T$, predictions are obtained by averaging the class probabilities $\mathbf{p}^i$ produced from each respective weight vector $\mathbf{w}[t]$ as specified in Eq.~\ref{eq:infer_prob}.

To enhance variation in rule extraction for the ensemble, the LLM is queried using a relatively high temperature. Furthermore, the sequence of few-shot exemplars is shuffled because prior work has demonstrated that the outcome of LLM queries can depend on this ordering~\cite{lu2022fantastically}\footnote{Further discussion regarding rule diversity is provided in Appendix~\ref{sec:diversity}}. Bagging is utilized to exploit predictive diversity, by selecting different subsets of features or data points for each trial, which becomes increasingly beneficial as the amount of training data or features grows. This ensemble approach introduces two notable strengths. First, mistakes in rule generation by the LLM in some rounds are compensated for by accurate rules from others, thereby bolstering the system's robustness. This aligns with the notion of self-consistency in LLMs~\cite{wang2022self}, since rules frequently generated across multiple rounds are usually trustworthy. Second, ensemble bagging mitigates the prompt length restrictions inherent to LLMs: for datasets whose tabular input surpasses the maximum allowable size, bagging helps reduce the required input, preserving model robustness. Though individual rule sets may fail to represent all feature relationships, assembling multiple weak learners across trials ensures that shortcomings in feature or instance coverage from single trials do not compromise overall predictive power.

\section{Experiments}
We assess the performance of \textsf{FeatLLM} on a range of tabular datasets and undertake various analytic procedures to elucidate the mechanics of our approach. Owing to space limitations, supplementary experimental results—including tests on different LLMs, qualitative investigations, and other analyses—are available in the Appendix.

\subsection{Performance Evaluation}
\paragraph{Datasets.}
Our experimental setup involves 13 datasets tailored for binary and multiclass classification tasks: (1) Adult~\cite{asuncion2007uci}, which involves predicting whether a person’s income exceeds \$50,000 per year; (2) Bank~\cite{moro2014data}, to determine whether a client subscribes to a term deposit; (3) Blood~\cite{yeh2009knowledge}, focused on forecasting whether a donor will return; (4) Car~\cite{kadra2021well}, aimed at assessing car quality; (5) Communities~\cite{misc_communities_and_crime_183}, for estimating regional crime rates; (6) Credit-g~\cite{kadra2021well}, which classifies individuals as posing either good or bad credit risks; (7) Diabetes\footnote{\texttt{kaggle.com/datasets/uciml/pima-indians-diabetes-database}}, concerned with predicting diabetes presence; (8) Heart\footnote{\texttt{kaggle.com/datasets/fedesoriano/heart-failure-prediction}}, which estimates coronary artery disease occurrences; (9) Myocardial~\cite{misc_myocardial_infarction_complications_579}, predicting chronic heart failure incidence.

The collection also encompasses newly introduced and synthetic datasets that are relatively underexplored and absent from LLMs’ pre-training data: (10) Cultivars, for forecasting the grain yield of soybean varieties\footnote{\texttt{archive.ics.uci.edu/dataset/913}}; (11) NHANES, for inferring age groups based on survey responses\footnote{\texttt{archive.ics.uci.edu/dataset/887}}; (12) Sequence-type, categorizing sequences as arithmetic, geometric, Fibonacci, or Collatz; (13) Solution-mix, discerning whether a mixture’s concentration from four solutions is above 0.5. The Cultivars and NHANES datasets were recently contributed to the UCI Machine Learning Repository (post September 2023), ensuring exclusion from the pre-training data of the utilized LLM API (gpt-3.5-turbo-0613). The remaining two datasets are synthetically constructed, precluding prior exposure or memorization by the LLM during its development.

A summary of dataset sizes, attributes, and intricacies can be found in Table~\ref{Tab:basic-desc}.
Each dataset includes explicit attribute naming and description. Notably, datasets like `Communities’ and `Myocardial’ feature over 100 attributes, posing notable context window limitations when using LLMs.

\vspace*{-0.12in}
\paragraph{Baselines.}
We benchmark against nine different baselines. The first three are established tabular data models: (1) Logistic Regression (LogReg), (2) XGBoost~\cite{chen2016xgboost}, and (3) RandomForest~\cite{ho1995random}. The following two—(4) SCARF~\cite{bahri2022scarf} and (5) STUNT~\cite{nam2023stunt}—are designed for semi-supervised learning with unlabeled data, though amassing ample unlabeled data may be impractical in real-world scenarios. Another contemporary baseline, (6) TabPFN~\cite{hollmann2022tabpfn}, leverages synthetic data generation across diverse distributions for pretraining. The final three methods employ LLMs: (7) In-context learning (In-context)~\cite{wei2022emergent}, (8) TABLET~\cite{slack2023tablet}, and (9) TabLLM~\cite{hegselmann2023tabllm}. In-context learning incorporates a handful of training samples directly into the LLM’s prompt without additional tuning. TABLET augments the input with external classifier-derived cues—specifically rule-sets and prototypes—to improve inference. TabLLM utilizes parameter-efficient training methods, such as IA3~\cite{liu2022few}, facilitating LLM adaptation with limited annotated examples.

\vspace*{-0.12in}
\paragraph{Implementation details.}
\textsf{FeatLLM} leverages GPT-3.5 as its foundational LLM but has been designed to operate independently of specific LLM architectures (refer to Appendix~\ref{sec:backbones} for comparative performance across multiple LLM backbones). The LLM inference process utilizes a temperature parameter of 0.5, with the top-p parameter remaining at the API's standard value of 1. We configured the ensemble count to 20 and extracted 10 rules, respectively. Insights into hyper-parameter sensitivity are provided in Figure~\ref{fig:hyperparameter2} as well as Appendix~\ref{sec:hyper-parameter}. For the linear modeling component, we optimized using Adam, setting the learning rate to 0.01 and training over 200 epochs. Epoch selection is supported through $k$-fold cross-validation. When replicating baseline methods, in-context learning is performed via GPT-3.5, while TabLLM is evaluated using T0~\cite{sanh2022multitask}, matching original experimental conditions. It is important to emphasize that TabLLM only permits LLMs with publicly released checkpoints; consequently, GPT-3.5 cannot be used within this framework. Standard machine learning algorithms—including LogReg, XGBoost, and RandomForest—are tuned through Grid-search in conjunction with $k$-fold cross-validation. The parameter $k$ is chosen to be 2 or 4, ensuring every class is represented at least once within the training split. Additional specifics regarding experimental setup are documented in Appendix~\ref{sec:baselines}.

\vspace*{-0.12in}
\paragraph{Results.}
Tables~\ref{tab:main1} and \ref{tab:main2} summarize the comparative evaluation across diverse datasets. For each dataset, we ran experiments three times, each with distinct random splits for training and testing, reporting the mean and standard deviation of AUC scores. The \textsf{FeatLLM} method persistently places either first or second relative to other baseline techniques. Remarkably, our model attains similar results with just 4 shots to those of standard tabular baselines trained with 32 shots (refer to Fig.~\ref{fig:compares}). Although models such as STUNT and TabPFN achieved notable gains in specific scenarios, their average advancement over classical machine learning approaches remained modest. In addition, LLM-powered models—including in-context learning, TABLET, and TabLLM—proved ineffective for tabular problems with extensive feature sets. By incorporating both bagging and ensembling, our method performs robustly on datasets with high-dimensional feature spaces. It is worth mentioning that the bagging and ensembling methodology we employed could, in principle, be integrated with other LLM-driven frameworks; nonetheless, such modifications would double the inference operations per sample, dramatically escalating computational costs and thus hindering practicality.

\subsection{Analysis \& Discussion}
\label{sec:analysis}
\paragraph{Ablation study.} We examine the impact of principal framework components by conducting ablation studies involving: (1) \textsf{-Tuning}: the exclusion of linear model weight adjustment, so that logits are computed simply as the sum of binary features; (2) \textsf{-Ensemble}: the removal of the ensembling step; (3) \textsf{-Description}: omitting feature descriptions; and (4) \textsf{-Reasoning}: bypassing the initial reasoning instruction during rule generation.

As reported in Table~\ref{tab:ablation}, all variations from the full method result in reduced AUC when averaged across datasets. Of note, the positive effect of weight tuning is most pronounced when the number of shots grows, implying that, with larger datasets, more precise rule weighting becomes impactful. In contrast, leveraging feature descriptions and reasoning guidance is particularly beneficial under few-shot conditions, underscoring the importance of harnessing LLMs' prior knowledge when data is limited. Across all experiment configurations, the use of ensembling persistently enhances performance.

\vspace*{-0.12in}
\paragraph{Training \& inference time comparison.} We evaluate and contrast the time required for training and inference across several approaches. All tests utilize the Adult dataset, incorporating a training set with 16 shots. The standard setup employs a single A100 GPU, except for TabLLM, which runs on four A100 GPUs due to model parallelism needs. Table~\ref{tab:cost} details these timing comparisons. Of particular note, \textsf{FeatLLM} achieves inference times that are relatively low and similar to traditional methods. Meanwhile, LLM-based alternatives—such as In-context learning, TABLET, and TabLLM—incur substantially greater inference time overhead. Regarding training duration, \textsf{FeatLLM} remains on par with other LLM-based models. It demands only 30 API calls, which do not substantially affect computational costs and can also be executed in parallel.

\vspace*{-0.12in}
\paragraph{Quality of features generated by \textsf{FeatLLM}.}
To assess the practical utility of rule-based features produced by \textsf{FeatLLM}, we conduct a straightforward evaluation. Specifically, we report the mean AUROC between each generated feature and the target labels, and then identify the proportion of features achieving an AUROC greater than 0.5 within every dataset. Table~\ref{tab:quality} summarizes our findings. The results demonstrate that most rules are strongly associated with the target attribute, thereby furnishing informative value for the task at hand (see the ``Before tuning" column). Additionally, when fine-tuning the linear model, features with marginal relevance—represented by weights below a predefined threshold—are automatically discarded (see the ``After tuning" column). The threshold applied is set at 0.1, based on analysis of the weight distribution.

\vspace*{-0.12in}
\paragraph{Effect of adding spurious correlations.} To evaluate the capacity of different approaches to mitigate the impact of noisy and spurious correlations under limited labeled data, we performed a targeted study. In this analysis, we utilized the Heart dataset, appending randomly selected columns from the Adult dataset that bear no relevance to the Heart prediction task. Model performance was assessed as the number of extraneous columns gradually increased. For a consistent benchmark, we used only 8 labeled instances and did not include an unlabeled sample pool, thereby omitting methods such as SCARF or STUNT from consideration.

Figure~\ref{fig:spurious} depicts the variation in results as spurious correlations are introduced. Notably, \textsf{FeatLLM} demonstrated superior resilience against performance deterioration stemming from noisy, irrelevant columns compared to the baseline methods. This outcome stems from the design of our approach, which integrates domain knowledge to inform feature-task relationships during rule extraction, minimizing reliance on irrelevant data and leading to greater robustness. Empirically, rules involving noisy columns constituted just 1-2\% of all derived rules. Nevertheless, it is also evident that \textsf{FeatLLM} is capable of learning spurious associations and may erroneously infer causality between features and targets if columns from datasets that are similar to the Heart dataset are arbitrarily added (refer to Appendix~\ref{sec:spurious-appendix}).

\vspace*{-0.12in}
\paragraph{Hyper-parameter analysis}
The functionality of \textsf{FeatLLM} hinges primarily on two hyper-parameters: the ensemble count and the quantity of rules per class extracted for each LLM prompt. To assess how these parameters affect model efficacy, we conducted a series of experimental studies. Our findings indicate that increasing the ensemble count generally boosts performance, though this improvement tapers off once the ensemble size approaches 20 (refer to Fig.~\ref{fig:appendix-hyperparameter1} in the Appendix). Regarding the optimal number of rules, although changes in this parameter do not yield statistically significant performance differences, our analysis suggests that setting this value to 10 strikes the best balance between keeping the prompt concise and maintaining strong predictive outcomes (see Fig.~\ref{fig:hyperparameter2}). We hypothesize that incorporating more rules expands the input feature space, which enhances information capacity but also elevates the risk of overfitting, particularly under the constraints of limited sample sizes.

\section{Conclusion}
We introduce \textsf{FeatLLM}, a method that elevates LLM-driven few-shot tabular learning to new heights. Unlike approaches that use LLMs simply to predict outcomes for individual samples, \textsf{FeatLLM} leverages LLMs as feature engineers, formulating explicit criteria through rule construction and aggregating them with bagging ensembles. This method significantly reduces inference costs, eliminates the need for model training beyond API interactions, and bypasses traditional limits imposed by feature dimensionality. As a result, \textsf{FeatLLM} not only delivers robust predictive performance but also emerges as a compelling, data-efficient alternative for LLM utilization in real-world tabular data contexts.

Notably, \textsf{FeatLLM} is tailored specifically for settings with scarce training data. In future work, we aim to broaden its scope to accommodate richer datasets, investigate additional feature forms beyond rules, and foster interpretability, thereby enabling the methodology’s application across diverse tasks and domains.

\section{Broader Impact}
\textsf{FeatLLM} is developed to exploit LLMs’ accumulated prior knowledge for superior performance in tabular learning, moving beyond conventional supervised learning bounds. By leveraging this prior knowledge, our framework tackles the persistent challenge of overfitting in scenarios where labeled data is limited, fostering greater data efficiency. This capability holds particular relevance for professionals in fields like finance and healthcare, where annotation is resource-intensive or hindered by data availability, and pretrained LLMs already encode substantial domain knowledge. A critical avenue for future exploration involves systematically examining how biases and inaccuracies embedded within LLM prior knowledge may influence, and possibly outweigh, the guidance provided by empirical labeled samples, especially in widespread and societally-sensitive deployments.

\end{document}